% =============================================================================
% Chapter 7 -- Debugging When Things Go Wrong
% =============================================================================
%
% Audience prerequisite: any reader who has at least attempted one run.
%
% Reader takeaway: can identify the cause of a failure from its symptoms via
%   the structured table, and knows the first two things to check before
%   opening an issue.
% =============================================================================

\chapter{Debugging When Things Go Wrong}
\label{ch:debug}

\refguide{10.4}{Tier I: Integration}

A \code{resolvent\_solve} call has a small number of ways to fail.
Each failure has a distinct symptom in the \code{ResolventResult} fields.
This chapter is a reference table: symptom, likely cause, and fix.

% =============================================================================
\section{Symptom--Cause--Fix Table}
\label{sec:debug-table}
% =============================================================================

The table below covers the ten most common failure modes.
Four of them (rows 1--4) are the canonical pitfalls from
Section~\ref{sec:pitfalls} in Chapter~\ref{ch:wrapping}.
Rows 5--10 are specific to GMRES behavior, conformal calibration, and
environment setup.

\medskip

\noindent
\begin{longtable}{p{0.28\textwidth} p{0.30\textwidth} p{0.33\textwidth}}
\toprule
\textbf{Symptom} & \textbf{Likely Cause} & \textbf{Fix} \\
\midrule
\endhead
\bottomrule
\endfoot

\raggedright\arraybackslash
GMRES converges (\texttt{residual\_history[-1] < gmres\_tol}) but
\texttt{delta\_x\_star} makes things worse when added to
\texttt{x\_star}.
&
\raggedright\arraybackslash
State-vector packing order mismatch.
The JVP unpacks or repacks the flat vector in depth-major order instead of
level-major order.
GMRES operates in the wrong coordinate system but the linear solve is
internally consistent, so convergence is unaffected.
&
\raggedright\arraybackslash
Run the round-trip test from \texttt{In[4]} in
Chapter~\ref{ch:wrapping}:
\texttt{assert torch.allclose(pops, x\_flat.reshape(L, N\_tau).T)}.
Fix the transposition convention in every JVP callback.
\\
\raggedright\arraybackslash
\texttt{torch.func.jvp} returns a zero tangent.
The host JVP is zero regardless of the input direction \texttt{v}.
&
\raggedright\arraybackslash
Autograd trace broken inside the function passed to
\texttt{torch.func.jvp}.
Common culprits: \texttt{.detach()} on the primal inside the function body;
an in-place operation (\texttt{tensor[i] = ...}) on a non-leaf tensor;
a NumPy round-trip that exits the computation graph.
&
\raggedright\arraybackslash
Remove all \texttt{.detach()} calls from inside the JVP function body.
Replace in-place operations with out-of-place equivalents.
Remove NumPy round-trips inside the traced function; convert to NumPy
only after returning from \texttt{torch.func.jvp}.
\\
\raggedright\arraybackslash
\texttt{RuntimeError} about mixed float32 and float64, or GMRES
converges far more slowly than expected (e.g., 100+ iterations at
moderate $\epsilon$).
&
\raggedright\arraybackslash
Float32 leaking from the APNN internal compute into the GMRES loop.
The APNN operates in float32 internally; float32 outputs that reach
GRAIL uncast silently lose eight decimal digits of GMRES precision.
&
\raggedright\arraybackslash
Add \texttt{v = v.to(dtype=torch.float64)} as the first line of every
JVP callback.
Cast \texttt{x\_star} and \texttt{delta\_T} explicitly before
passing them to \texttt{resolvent\_solve}.
See the advisory box in Section~\ref{pit:dtype}.
\\
\raggedright\arraybackslash
\texttt{delta\_x\_star} is near zero regardless of the input
\texttt{delta\_T}, or has the right magnitude but the wrong sign.
&
\raggedright\arraybackslash
\texttt{emulator\_jvp} uses the APNN's own autograd Jacobian instead of
the truth SE Jacobian.
An MSE-trained APNN matches truth function values to $\sim\!10^{-5}$
relative, but its derivative can differ by $10^5\times$ because MSE
training does not supervise derivatives.
Using the APNN gradient in the GMRES operator gives a wrong answer.
&
\raggedright\arraybackslash
\texttt{emulator\_jvp} must differentiate the \emph{truth} SE solver
(e.g., \texttt{rate\_equation\_solve}), not the APNN.
The APNN enters GRAIL only through \texttt{delta\_T}, the per-call
residual.
See Section~\ref{pit:jacobian} and \texttt{In[19]} in
Chapter~\ref{ch:wrapping}.
\\
\raggedright\arraybackslash
GMRES residual history is monotone decreasing but stagnates well above
\texttt{gmres\_tol} (e.g., plateaus at $\sim 10^{-2}$--$10^{-3}$ and
does not further decrease even at \texttt{gmres\_maxiter}).
&
\raggedright\arraybackslash
Spectral radius of $D\Tcyc$ is at or above 1.
This occurs at deep $\epsilon$ ($\lesssim 10^{-5}$) and large
$\tau_\mathrm{max}$ ($\gtrsim 10^4$), where the cycle Jacobian has
eigenvalues approaching 1 from below.
The Krylov subspace cannot reduce the residual further without a
preconditioner.
&
\raggedright\arraybackslash
Add a $\Lambda^*$-style diagonal preconditioner to \texttt{resolvent\_solve}
via the \texttt{preconditioner} argument.
See Section~\ref{sec:precond} in Chapter~\ref{ch:perf} for the
construction.
As a quick check, verify the condition estimate:
if \texttt{condition\_estimate > 1e3}, a preconditioner is warranted.
\\
\raggedright\arraybackslash
GMRES residual grows after some initial decrease
(\texttt{residual\_history} is non-monotone, with the final value larger
than the minimum).
&
\raggedright\arraybackslash
Floating-point instability in the GMRES update.
The most common cause is a dtype inconsistency: float32 leaking into
a float64 computation creates catastrophic cancellation in the Arnoldi
orthogonalization.
A secondary cause is a state-vector layout bug that produces a
non-Hermitian effective operator.
&
\raggedright\arraybackslash
Check dtype first (Pitfall~\ref{pit:dtype}).
Then run the round-trip test (Pitfall~\ref{pit:layout}).
Then check the JVP finite-difference consistency (Step~2 of the
debugging recipe in Section~\ref{sec:debug-recipe} in
Chapter~\ref{ch:wrapping}).
\\
\raggedright\arraybackslash
\texttt{condition\_estimate} is \texttt{inf}.
&
\raggedright\arraybackslash
This is not an error.
\texttt{inf} means the Arnoldi Hessenberg has a zero singular value,
which occurs when GMRES finds an exact solution in the current Krylov
subspace.
&
\raggedright\arraybackslash
No action needed.
An \texttt{inf} condition estimate with a final residual below
\texttt{gmres\_tol} is a healthy sign of fast convergence.
\\
\raggedright\arraybackslash
\texttt{condition\_estimate > 1e12} with the final residual stuck above
\texttt{gmres\_tol}.
&
\raggedright\arraybackslash
The cycle's spectral radius is pinned at 1.
At extreme $\epsilon$ ($\lesssim 10^{-6}$) the operator
$(I - D\Tcyc)$ is nearly singular.
&
\raggedright\arraybackslash
Add a preconditioner (Chapter~\ref{ch:perf} Section~\ref{sec:precond}).
Alternatively, reduce \texttt{gmres\_tol} to a value that GMRES can
actually reach, and accept that the correction is first-order only.
Check that the atmosphere's $\tau_\mathrm{max}$ is inside the linearity
validity regime.
\\
\raggedright\arraybackslash
\texttt{||delta\_T||} is suspiciously small ($< 10^{-12}$).
&
\raggedright\arraybackslash
\texttt{delta\_T} is computed against training data, not a held-out
validation point.
The APNN memorizes training inputs to machine precision, so
$\dT$ at a training point is effectively zero.
A near-zero \texttt{delta\_T} produces a near-zero \texttt{delta\_x\_star}
regardless of the true amplification.
&
\raggedright\arraybackslash
Compute \texttt{delta\_T = T\_emul(J\_bar) - T\_truth(J\_bar)} using a
fresh validation $\bar{J}$ not seen during training.
The standard practice is to generate 20\% of the training configs as a
held-out validation set and use only those $\bar{J}$ values when
computing \texttt{delta\_T}.
\\
\raggedright\arraybackslash
\texttt{pytest} exits at collection with
\texttt{OMP: Error \#15: Initializing libiomp5.dylib, but found
libiomp5.dylib already initialized.}
&
\raggedright\arraybackslash
PyTorch and h5py ship conflicting \texttt{libomp.dylib} copies in the
conda environment.
The dynamic linker initializes the library twice and aborts.
&
\raggedright\arraybackslash
Set \texttt{KMP\_DUPLICATE\_LIB\_OK=TRUE} in the environment before
launching pytest:
\begin{lstlisting}[style=grailInStyle]
KMP_DUPLICATE_LIB_OK=TRUE \
    python -m pytest -m integration
\end{lstlisting}
The workaround is already applied programmatically in
\texttt{tests/conftest.py} for test execution.
The production fix is to align both packages against the same
\texttt{libomp} at install time.
\\

\bottomrule
\end{longtable}

% =============================================================================
\section{Reproducing a Failure Locally}
\label{sec:repro-locally}
% =============================================================================

When a failure is intermittent or depends on a large ensemble, the fastest
path to a fix is to minimize the reproducer to a single-config standalone
script.
A systematic binary-search halving works in all cases.

The idea: run the full failing sweep, identify which half fails, discard
the passing half, repeat until a single config remains.

\grailInLabel{1}
\begin{lstlisting}[style=grailInStyle]
# Binary-search failure isolation on a list of configs.
# 'configs' is a list of dicts with the parameters for one resolvent_solve call.
# 'run_batch' returns a list of booleans (True = pass, False = fail).

def bisect_failure(configs, run_batch, depth=0, max_depth=10):
    """Returns the first failing config, or None if all pass."""
    if len(configs) == 0 or depth >= max_depth:
        return None
    results = run_batch(configs)
    if all(results):
        return None   # All pass -- no failure in this batch.
    if len(configs) == 1:
        return configs[0]   # Isolated to one config.

    mid = len(configs) // 2
    left, right = configs[:mid], configs[mid:]
    left_results, right_results = results[:mid], results[mid:]

    if not all(left_results):
        return bisect_failure(left, run_batch, depth + 1)
    else:
        return bisect_failure(right, run_batch, depth + 1)

# Example: 16 configs from the epsilon sweep.
epsilon_configs = [
    {"epsilon": 10.0 ** (-k), "tau_max": 1e4, "n_tau": 40}
    for k in range(2, 7)
]

def run_batch_stub(cfgs):
    # Replace with your actual resolvent_solve call that returns True/False.
    return [cfg["epsilon"] > 1e-5 for cfg in cfgs]

failing = bisect_failure(epsilon_configs, run_batch_stub)
print(f"Failing config isolated: {failing}")
\end{lstlisting}

\grailOutLabel{1}
\begin{lstlisting}[style=grailOutStyle]
Failing config isolated: {'epsilon': 1e-05, 'tau_max': 10000.0, 'n_tau': 40}
\end{lstlisting}

\noindent With the failing config isolated, build a standalone script that
reproduces the failure in under 30 seconds:

\grailInLabel{2}
\begin{lstlisting}[style=grailInStyle]
# Standalone failure reproducer template.
# Replace the body with the actual failing config parameters.

import os
os.environ.setdefault("KMP_DUPLICATE_LIB_OK", "TRUE")

import torch
from grail import resolvent_solve

# --- Config from binary search ---
epsilon = 1e-5
tau_max = 1e4
n_tau   = 40

# --- Minimal atmosphere ---
from atmosphere.atmosphere import Atmosphere
from atmosphere.grid import DepthGrid, FrequencyGrid, AngleQuadrature

depth     = DepthGrid.log_spaced(tau_min=1e-4, tau_max=tau_max, n_tau=n_tau)
frequency = FrequencyGrid.doppler(n_nu=11, half_width=5.0)
angle     = AngleQuadrature.gauss_legendre(n_mu=3)
atm = Atmosphere.isothermal(
    T0=6000.0, n_e0=1e10, epsilon=epsilon, nu_0=2.47e15,
    depth=depth, frequency=frequency, angle=angle,
)

# --- Minimal JVPs (copy-paste from your actual integration) ---
# host_jvp = ...
# emul_jvp = ...
# x_star   = ...
# delta_T  = ...

# --- Call and inspect ---
# result = resolvent_solve(x_star, delta_T, emul_jvp, host_jvp)
# print(result.residual_history)
# print(result.condition_estimate)
print("Reproducer built. Plug in your JVPs and uncomment the call.")
\end{lstlisting}

\grailOutLabel{2}
\begin{lstlisting}[style=grailOutStyle]
Reproducer built. Plug in your JVPs and uncomment the call.
\end{lstlisting}

\grailNote{The standalone reproducer is the required minimum for a GitHub
issue.
Paste it into the issue body verbatim.
If the reproducer takes longer than 30 seconds to reach the failure,
halve \texttt{n\_tau} and \texttt{n\_nu} until it does.
A reproducer that takes five minutes discourages fast turnaround.}

% =============================================================================
\section{Filing a GitHub Issue}
\label{sec:github-issue}
% =============================================================================

When the table in Section~\ref{sec:debug-table} and the binary-search
reproducer in Section~\ref{sec:repro-locally} do not resolve the failure,
file a GitHub issue.
Include the following.

\begin{enumerate}
  \item \textbf{The standalone reproducer} from Section~\ref{sec:repro-locally}.
        Under 100 lines, runs in under 30 seconds, reproduces the failure without
        external data files.

  \item \textbf{The full \texttt{ResolventResult} field dump.}
        Print every field:
\grailInLabel{3}
\begin{lstlisting}[style=grailInStyle]
print(f"iter_count:         {result.iter_count}")
print(f"residual_history:   {result.residual_history}")
print(f"condition_estimate: {result.condition_estimate:.4e}")
print(f"pca_rank:           {result.pca_rank}")
print(f"delta_x_star norm:  {result.delta_x_star.norm():.4e}")
print(f"delta_T norm:       {delta_T.norm():.4e}")
print(f"x_star norm:        {x_star.norm():.4e}")
\end{lstlisting}

  \item \textbf{The GMRES residual history as a list.}
        Copy the raw list, not a plot.
        The shape of the curve (monotone decrease, stagnation plateau, growth)
        determines the diagnostic category.

  \item \textbf{The config parameters.}
        At minimum: $\epsilon$, $\tau_\mathrm{max}$, $N_\tau$, $N_\nu$, $N_\mu$,
        $L$ (number of atomic levels), and the APNN architecture (input/output
        dimension, number of layers).

  \item \textbf{The APNN checkpoint hash.}
        If a trained checkpoint is involved, include
        \texttt{sha256sum examples/data/structured/model.pt} or equivalent.
        Checkpoint hashes let the maintainer reproduce the exact model.
\grailInLabel{4}
\begin{lstlisting}[style=grailInStyle]
import hashlib, pathlib
ckpt_path = pathlib.Path("examples/data/structured/model.pt")
if ckpt_path.exists():
    digest = hashlib.sha256(ckpt_path.read_bytes()).hexdigest()[:16]
    print(f"APNN checkpoint sha256 prefix: {digest}")
else:
    print("No checkpoint found at examples/data/structured/model.pt")
\end{lstlisting}

  \item \textbf{The environment.}
        Output of:
\grailInLabel{5}
\begin{lstlisting}[style=grailInStyle]
import torch, numpy, sys
print(f"Python:  {sys.version}")
print(f"PyTorch: {torch.__version__}")
print(f"NumPy:   {numpy.__version__}")
\end{lstlisting}

  \item \textbf{Which row in Section~\ref{sec:debug-table} best matches the symptom}
        and why the suggested fix did not resolve it.
        This narrows the triage immediately.

\end{enumerate}

\grailNote{Issues filed without a standalone reproducer are deprioritized.
The maintainers cannot reproduce a failure from a description alone.
A 50-line script that runs in 10 seconds is the single most helpful thing
you can include.}
